\documentclass[10pt,a4paper]{amsart}
\usepackage[margin=19mm]{geometry}
\usepackage{amsmath,amssymb,mathtools}
\usepackage[scr=rsfs,cal=euler]{mathalfa}
\usepackage{xfrac}
\usepackage[unicode,hidelinks,bookmarks=false]{hyperref}
\newcommand{\ie}{i.e.\ }
\newcommand{\cf}{cf.\ }
\newcommand{\resp}{respectively\ }
\newcommand{\Subsection}{\subsection}
\providecommand{\nobreakdash}{\nobreak\hbox{-}}
\setlength{\emergencystretch}{2em}
\multlinegap0pt
\newcommand{\K}{\mathbb{K}}
\usepackage{amssymb}
\usepackage[shortlabels]{enumitem}
\newtheorem{theorem}{Theorem}
\title{\bf Prime numbers and factorization of polynomials}
\author{Jitender Singh}
\date{Source: arXiv:2411.18366v1 (CC BY 4.0)}
\begin{document}
\maketitle

\noindent\emph{Source and licence.} \bf Prime numbers and factorization of polynomials. Version arXiv:2411.18366v1 (CC BY 4.0), distributed by arXiv under the Creative Commons Attribution 4.0 International licence, http://creativecommons.org/licenses/by/4.0/, as recorded in the arXiv licence metadata for this exact version and shown on the abstract page https://arxiv.org/abs/2411.18366v1. Full licence notice: \texttt{licenses/arxiv-2411.18366-CC-BY-4.0.txt}.\par\smallskip

\noindent\textit{Editorial scope.} Counted unit: the article's theorem th4 (source0.tex lines 201--211). The article's complete original proof of it is delivered with it (source0.tex lines 285--304, one \texttt{proof} environment of 20 lines, which the article places away from the statement). No mathematics is written, repaired or re-derived: the statement and the proof are the article's own lines, in the article's own notation, and no retained line is substituted or re-flowed.  No further source-local context is needed: every cross-reference of every retained line resolves inside the two delivered spans.  Nothing was omitted from any retained span, so the package carries no omission record.  No retained line contains a citation, so the wrapper carries no bibliography and no bibliography entry.  No retained line contains a figure environment, an image-inclusion command or drawing code, so no image is delivered and the package declares no asset dependency.  Editorial formatting only, in addition: an AMS \texttt{amsart} preamble that restates the article's own declarations and macros used by the retained lines, namely the environments \texttt{theorem} on separate counters, together with the standard packages those lines need.  The retained environments print in this document as Theorem 1 in order of appearance, whereas the article prints the same items with the numbering of its own full document.  The complete native source of the article as distributed in the arXiv e-print for this exact version is bundled unchanged as \texttt{sources/169225/source0.tex}.
\begin{theorem}\label{th4}
 Let $\K$ be a field, and let $f=a_0(x)+a_1(x)y+\cdots+a_n(x)y^n\in \K[x,y]$ be a  polynomial with $a_0a_n\neq 0$, and for a real number $\rho>1$, define
 \begin{eqnarray*}
 H_f=\rho^{\max_{0\leq i\leq n-1}\deg a_i-\deg a_n}.
 \end{eqnarray*}
  Suppose there exists a polynomial $a\in \K[x]$ for which
 \begin{eqnarray*}
              \deg a\geq  \frac{\log(H_f+2)}{\log \rho},
 \end{eqnarray*}
  and $\nu_a$ is the number of irreducible factors of the polynomial $f(x,a(x))$ in $\K[x]$ counted with multiplicities. Then  the polynomial $f(x,y)$ is a product of at most $\nu_a$ irreducible factors in $\K[x,y]$. In particular, if $\nu_a=1$, then $f$ is irreducible in $\K[x,y]$.
\end{theorem}

\begin{proof}[\bf Proof of Theorem \ref{th4}] First note that $\deg a\geq  \log (H_f+2)/\log \rho$ if and only if $\|a\|\geq H_f+2$.  For any $y\in \overline{\K(x)}$ satisfying $\|y\|\geq H_f+1$, we have on using properties of $\|\cdot\|$ that
\begin{eqnarray*}
\frac{\|f(x,y)\|}{\|y^n a_n\|} \geq 1-\sum_{i=1}^{n-1}\frac{\|a_i\|}{\|a_n\|}\frac{1}{\|y^{n-i}\|}
 &= & 1-\sum_{i=1}^{n-1}\rho^{\deg a_i-\deg a_n}\frac{1}{\|y\|^{n-i}}\\
 &\geq & 1-\rho^{\max_{i}\deg a_i-\deg a_n}\sum_{i=1}^{n-1}\frac{1}{H_f^{n-i}}\\
   &\geq & 1-H_f\sum_{i=1}^{n-1}\frac{1}{H_f^{n-i}}= \bigl(H_f+1\bigr)^{-n}>0.
\end{eqnarray*}
This shows that any zero $\theta\in \overline{\K(x)}$ of $f(x,y)=a_0+a_1 y+\cdots+a_ny^n\in \K[x,y]$ satisfies $\|\theta\|<H_f+1$.

Now assume that $f(x,y)=f_1(x,y)\cdots f_r(x,y)$ be a product of $r$ irreducible polynomials $f_i$'s in $\K[x,y]$. Then $\|f_i(x,a(x))\|\geq 1$ for every $i=1,\ldots,r$. If $\alpha_i(x)\in \K[x]$ is the leading coefficient of $f_i$, then we may write
\begin{eqnarray*}
f_i(x,y) &=& \alpha_i(x)\prod_{\theta} (y-\theta),
\end{eqnarray*}
where the product runs over all zeros of $f_i$ in $\overline{\K(x)}$. We then have
\begin{eqnarray*}
\|f_i(x,a(x))\| = \|\alpha_i(x)\|\prod_{\theta} \|a(x)-\theta\|&\geq& \prod_{\theta} (\|a(x)\|-\|\theta\|)\\
&>& \prod_{\theta} (H_f+2-(H_f+1))=1,
\end{eqnarray*}
which shows that $\deg f_i(x,a(x))>0$ for each $i$. Consequently, it follows from the equality $f(x,a(x))=f_1(x,a(x))\cdots f_r(x,a(x))$ that there must be an irreducible factor of $f(x,a(x))$  dividing $f_i(x,a(x))$ for each $i$. This proves that $r\leq \nu_a$.
\end{proof}

\end{document}
